\subsection{可解圈空间的万有覆叠}

设 B 是一个拓扑空间，b 是 B 的一点。回忆 \(\Lambda_{b}(\mathrm{B})\) 表示 \(\mathcal{C}_{c}(\mathbf{I};\mathrm{B})\) 中由以 b 为起点的道路组成的子空间，赋以紧收敛拓扑的商拓扑。我们用 \(e_{\mathrm{B}}\colon\Lambda_{b}(\mathrm{B})\to\mathrm{B}\) 表示把一条道路对应到其终点的映射；它是连续的。用 \(\lambda_{b}(\mathrm{B})\) 表示 \(\Lambda_{b}(\mathrm{B})\) 关于严格同伦关系的商空间，并赋以商拓扑。由于两条严格同伦的道路有相同的终点，映射 \(e_{B}\) 通过过渡到商，定义一个连续映射 \(\varepsilon_{\mathrm{B}}\colon\lambda_{b}(\mathrm{B})\to\mathrm{B}\)，称为终点映射。

定理 1. --- 设 B 是一个连通且局部道路连通的拓扑空间，并设 b 是 B 的一点。下列性质等价：

\begin{enumerate}
\def\labelenumi{(\roman{enumi})}
\item
  空间 B 是可解圈的。
\item
  存在 B 的一个非空且道路单连通的覆叠。
\item
  空间 \(\lambda_{b}(\mathrm{B})\)，赋以终点映射 \(\varepsilon_{\mathrm{B}}\colon\lambda_{b}(\mathrm{B})\to\mathrm{B}\)，是 B 的一个覆叠。
\item
  群 \(\pi_1(B, b)\) 对容许拓扑是离散的。
\item
  群 \(\pi_{1}(B,b)\) 对紧收敛拓扑的商拓扑是离散的。
\end{enumerate}

此外，当这些条件满足时，空间 \(\lambda_{b}(\mathrm{B})\) 是道路单连通的、单连通的、以 \(\pi_{1}(\mathrm{B},b)^{\circ}\) 为群的伽罗瓦覆叠，且点覆叠 \((\lambda_{b}(\mathrm{B}),\varepsilon_{b})\) 是点空间 \((\mathrm{B},b)\) 的一个万有覆叠。

由可解圈空间的定义可知，\(\pi_{1}(B,b)\) 中退缩为单位元的子群是容许的，当且仅当 B 是可解圈的。这就证明了 (i)⇔(iv)。此外，性质 (iv) 与 (v) 的等价性由 III, p. 320 的注 4 与注 5 得出。

(iv)⇒(ii)。由 III, p. 316, 命题 5，存在 B 的一个覆叠 (E, p) 及一点 \(x \in E_{b}\)，使 \(p_{*}(\pi_{1}(E, x)) = \{e\}\)；特别地，\(\pi_{1}(E, x)\) 退缩为单位元。于是 x 在 E 中的弧连通分支是 B 的一个非空覆叠（I, p. 80, 命题 6 的推论 1），且是道路单连通的。

(ii)⇒(iii)。设 E 是 B 的一个非空道路单连通覆叠。设 x 是 \(E_{b}\) 的一点（这样的点存在，因为 B 连通，I, p. 74, 命题 4）。投影 \(q: E \to B\) 诱导一个同胚 \(\Lambda_{x}(E) \to \Lambda_{b}(B)\)（III, p. 302, 命题 3 的推论 2），从而通过过渡到弧连通分支，得到一个同胚 \(q_{*}: \lambda_{x}(E) \simeq \lambda_{b}(B)\)。终点映射 \(\Lambda_{x}(E) \to E\) 是连续且开的，因为 E 局部道路连通（III, p. 264, 推论）。因此它通过过渡到弧连通分支所定义的映射 \(\varepsilon_{\mathrm{E}}: \lambda_{x}(E) \to E\) 是连续且开的。映射 \(\varepsilon_{\mathrm{E}} \circ (q_{*})^{-1}: \lambda_{b}(B) \to E\) 是双射的、连续的且开的。这就证明了赋以紧收敛拓扑与终点映射的拓扑空间 \(\lambda_{b}(B)\) 是 B 的一个覆叠。

(iii)⇒(v)。事实上，赋以紧收敛拓扑的商拓扑的集合 \(\pi_{1}(B,b)\)，与 B-空间 \(\lambda_{b}(B)\) 在 b 上、在 b 处常值圈所属的那个类处的纤维等同。若 \(\lambda_{b}(B)\) 是 B 的覆叠，则 \(\pi_{1}(B,b)\) 对紧收敛拓扑是离散的。

假设这些断言成立。由前文，空间 \(\lambda_b(B)\) 是 B 的一个道路单连通的覆叠，因而是单连通的。于是点空间 ( \(\lambda_b(B), \varepsilon_b\) ) 是 (B, \(b\) ) 的一个万有覆叠（I, p. 126, 命题 3 的推论），且覆叠 \(\lambda_b(B)\) 是 B 的一个伽罗瓦覆叠（I, p. 120, 命题 1）。典范同态 \(h_{(\lambda_b(B), \varepsilon_b)}: \pi_1(B, b) \to \text{Aut}_B(\lambda_b(B))\) 于是是一个同构（III, p. 306, 命题 5）。

注 1. --- 设 B 是一个连通可解圈空间，并设 b 是 B 的一点。由定理 1, (iv) 可知，群 \(\pi_{1}(B,b)\) 的每个作用都是容许的。因此，每个 \(\pi_{1}(B,b)\) -集合都同构于一个 \(\pi_{1}(B,b)\) -集合 \(E_{b}\)，其中 E 是 B 的覆叠（III, p. 318, 命题 7）。还回忆（III, p. 310, 命题 2）：若 E 与 \(E'\) 是 B 的覆叠，则映射 \(f \mapsto f_{b}\) 诱导一个从 E 到 \(E'\) 中的 B-空间态射所成之集到 \(\pi_{1}(B,b)\) -集合的从 \(E_{b}\) 到 \(E'_{b}\) 上的态射所成之集上的双射。

用范畴的语言来说：把 B 的覆叠 E 对应到纤维 \(E_b\) 的函子，是从 B 的覆叠的范畴到 \(\pi_1(B,b)\) -集合的范畴的一个范畴等价。*

推论 1. --- 设 B 是一个可解圈拓扑空间，并设 b 是 B 的一点。设 E 是 B 的一个连通覆叠，x 是纤维 \(E_{b}\) 的一点。下列性质等价：

\begin{enumerate}
\def\labelenumi{(\roman{enumi})}
\item
  群 \(\pi_{1}(\mathrm{E},x)\) 是平凡的。
\item
  空间 E 是单连通的。
\item
  点空间 (E, x) 是 (B, b) 的一个万有覆叠。
\end{enumerate}

蕴涵关系 (i)⇒(ii) 已仅在“空间 B 连通且局部道路连通”这一假设下证明（III, p. 309, 命题 1 的推论 2），而蕴涵关系 (ii)⇒(iii) 则在没有任何假设下证明（I, p. 126, 推论）。

我们来证明 (iii)⇒(i)。由 IV, p. 342 的定理 1，存在 B 的一个覆叠 E' 及一点 \(x'\)（属于纤维 \(E'_{b}\)），使群 \(\pi_1(E', x')\) 退缩为单位元。在假设 (iii) 下，存在

一个点态 B-态射 \(f \colon (\mathrm{E}, x) \to (\mathrm{E}', x')\)。用 \(p \colon \mathrm{E} \to \mathrm{B}\) 与 \(p' \colon \mathrm{E}' \to \mathrm{B}\) 表示覆叠 \(\mathrm{E}\) 与 \(\mathrm{E}'\) 的投影；我们有 \(p = p' \circ f\)，因此 \(\pi_1(p, x) = \pi_1(p', x') \circ \pi_1(f, x)\)。由于群 \(\pi_1(\mathrm{E}', x')\) 平凡，单射同态 \(\pi_1(p, x)\) 以单位元为其像。这就证明了群 \(\pi_1(\mathrm{E}, x)\) 退缩为单位元。

推论 2. --- 设 B 是一个可解圈拓扑空间，并设 b 是 B 的一点。设 (E, x) 是点空间 (B, b) 的一个万有覆叠。对 B 的每个覆叠 E' 及每一点 \(x' \in E'_{b}\)，使 f(x) = x' 的唯一的 B-态射 f: E → E' 使 E 成为 E' 的一个覆叠。于是点空间 (E, x) 是 (E', x') 的一个万有覆叠。

由 I, p. 81 的命题 7，赋以映射 f 的空间 E 是 E' 的一个覆叠。最后的断言于是由推论 1 得出。

设 B 是一个可解圈拓扑空间，并设 \(a\)、\(b\) 是 B 的两点。商拓扑空间 \(\varpi_{a,b}(\mathrm{B})\) 同胚于赋以紧收敛拓扑的商拓扑的空间 \(\pi_1(\mathrm{B},a)\)（III, p. 293, 注 3），因而由 IV, p. 342 的定理 1 是离散的。所以，B 中连接 \(a\) 与 \(b\) 的每条道路的严格同伦类都是 \(\Lambda_{a,b}(\mathrm{B})\) 的一个开子集。更一般地：

命题 4. --- 设 B 是一个可解圈拓扑空间。严格同伦是拓扑空间 \(\mathcal{C}_{c}(\mathbf{I};\mathbf{B})\) 中的一个开等价关系。

首先假设空间 B 是道路单连通的。设 \(\varphi\colon\Lambda(B)\to B\times B\) 是把 B 中的道路 \(c\) 对应到由其起点与终点组成的偶 \((c(0),c(1))\) 的映射。为使 B 中的两条道路严格同伦，必要且充分的是它们有相同的起点与相同的终点（III, p. 292, 命题 2）。由于空间 B 连通且局部道路连通，映射 \(\varphi\) 是满射的、连续的且开的（III, p. 262, 命题 10）。因此严格同伦关系是开的（TG, I, p. 32, 命题 3）。

在一般情形，设 E 是 B 的一个道路单连通覆叠，当 \(B \neq \varnothing\) 时非空。由于 B 连通，这个覆叠是满射的（I, p. 74, 命题 4）。用 p 表示 B-空间 E 的投影，并用 \(\widetilde{p} \colon \Lambda(\mathrm{E}) \to \Lambda(\mathrm{B})\) 表示把 E 中的道路 c 对应到道路 \(p \circ c\) 的映射。赋以映射 \(\widetilde{p}\) 的 \(\Lambda(\mathrm{E})\) 是 \(\Lambda(\mathrm{B})\) 的一个覆叠（III, p. 302, 命题 3 的推论 1）。为使 B 中的两条道路严格同伦，必要且充分的是它们都是 E 中两条严格同伦的道路在映射 \(\widetilde{p}\) 下的像（III, p. 302, 命题 4）。设 U 是空间 \(\Lambda(B)\) 的一个开子集。集合 \((\widetilde{p})(U)\) 在 \(\Lambda(E)\) 中开；由证明的第一部分，集合 \(U'\)（关于严格同伦关系、相对于 \((\widetilde{p})(U)\) 饱和）在 \(\Lambda(E)\) 中开。由于 \(\Lambda(E)\) 是 \(\Lambda(B)\) 的覆叠，集合 \(\widetilde{p}(U')\) 在 \(\Lambda(B)\) 中开。由于 \(\widetilde{p}(U')\) 是 U 关于严格同伦关系的饱和化，命题得证。

